% Original source: Zhou, physical PDF pp. 158--162 (printed pp. 142--146).
% Figures: none.
% Uncertainty log: footnote 7 begins outside this assigned page range.

% ZHOU-SOURCE-PAGE: 158 PRINTED: 142

\begin{solution}
In the last problem, we mentioned that the payoff of a put is a convex function in
stock price. The price of a put option as a function of the strike price is a convex
function as well. Since a put option with strike 0 is worthless, we always have

\[
P(0)+\lambda P(K)=\lambda P(K)>P(\lambda K).
\]

For this specific problem, we should have $8/9\cdot P(90)=8/9\cdot9=8>P(80)$. Since the
put option with strike price \$80 is currently priced at 8, it is overpriced and we should
short it. The overall arbitrage portfolio is to short 9 units of put with $K=\$80$ and long
8 units of put with $K=90$. At time 0, the initial cash flow is 0. At the maturity date, we
have three possible scenarios:

\[
S_T\ge 90,\qquad \text{payoff}=0\quad\text{(No put is exercised.)}
\]

\[
90>S_T\ge 80,\qquad \text{payoff}=8\cdot(90-S_T)>0\quad\text{(Puts with }K=90\text{ are exercised.)}
\]

\[
S_T<80,\qquad \text{payoff}=8\cdot(90-S_T)-9\cdot(80-S_T)=S_T>0\quad\text{(All puts are exercised.)}
\]

The final payoff $\ge0$ with positive probability that payoff $>0$. So it is clearly an
arbitrage opportunity.

\end{solution}

\subsubsection*{Black-Scholes-Merton differential equation}

\begin{problem}{Can you write down the Black-Scholes-Merton differential equation and briefly explain
how to derive it?}
\end{problem}

\begin{solution}
If the evolution of the stock price is a geometric Brownian motion,
$\dd S=\mu S\,\dd t+\sigma S\,\dd W(t)$, and the derivative $V=V(S,t)$ is a function of $S$ and $t$,
then applying Ito's lemma yields:

\[
\begin{aligned}
\dd V&=\left(\frac{\partial V}{\partial t}+\mu S\frac{\partial V}{\partial S}
+\frac{1}{2}\sigma^2S^2\frac{\partial^2V}{\partial S^2}\right)\,\dd t\\
&+\sigma S\frac{\partial V}{\partial S}\,\dd W(t),\\
&\quad\text{where }W(t)\text{ is a Brownian motion.}
\end{aligned}
\]

The Black-Scholes-Merton differential equation is a partial differential equation that
should be satisfied by $V$:

\[
\frac{\partial V}{\partial t}+rS\frac{\partial V}{\partial S}
+\frac{1}{2}\sigma^2S^2\frac{\partial^2V}{\partial S^2}=rV.
\]

To derive the Black-Scholes-Merton differential equation, we build a portfolio with two
components: long one unit of the derivative and short $\dfrac{\partial V}{\partial S}$ unit
of the underlying stock.

Then the portfolio has value $\Pi=V-\dfrac{\partial V}{\partial S}S$ and the change of
$\Pi$ follows equation

% ZHOU-SOURCE-PAGE: 159 PRINTED: 143

\[
\begin{aligned}
\dd\Pi&=\dd V-\frac{\partial V}{\partial S}\,\dd S\\
&=\left(\frac{\partial V}{\partial t}+\mu S\frac{\partial V}{\partial S}
+\frac{1}{2}\sigma^2S^2\frac{\partial^2V}{\partial S^2}\right)\,\dd t\\
&\quad+\sigma S\frac{\partial V}{\partial S}\,\dd W(t)\\
&\quad-\frac{\partial V}{\partial S}\left(\mu S\,\dd t+\sigma S\,\dd W(t)\right)\\
&=\left(\frac{\partial V}{\partial t}+\frac{1}{2}\sigma^2S^2\frac{\partial^2V}{\partial S^2}\right)\,\dd t.
\end{aligned}
\]

It is apparent that this portfolio is risk-free since it has no diffusion term. It should
have risk-free rate of return as well: $\dd\Pi=r\left(V-\dfrac{\partial V}{\partial S}S\right)\,\dd t$.
Combining these results we have

\[
\begin{aligned}
\left(\frac{\partial V}{\partial t}+\frac{1}{2}\sigma^2S^2\frac{\partial^2V}{\partial S^2}\right)\,\dd t
&=r\left(V-\frac{\partial V}{\partial S}S\right)\,\dd t\\
&\Rightarrow \frac{\partial V}{\partial t}+rS\frac{\partial V}{\partial S}\\
&\quad+\frac{1}{2}\sigma^2S^2\frac{\partial^2V}{\partial S^2}=rV,
\end{aligned}
\]

which is the Black-Scholes-Merton differential equation.

The Black-Scholes-Merton differential equation is a special case of the discounted
Feynman-Kac theorem. The discounted Feynman-Kac theorem builds the bridge between
stochastic differential equations and partial differential equations and applies to all Ito
processes in general:

Let $X$ be an Ito process given by equation
$\dd X(t)=\beta(t,X)\,\dd t+\gamma(t,X)\,\dd W(t)$ and $f(x)$ be a function of $X$. Define function
$V(t,x)=\E\left[e^{-r(T-t)}f(X_T)\vert X_t=x\right]$, then $V(t,x)$ is a martingale process
that satisfies the partial differential equation

\[
\begin{aligned}
\frac{\partial V}{\partial t}+\beta(t,x)\frac{\partial V}{\partial x}\\
&\quad+\frac{1}{2}\gamma^2(t,x)\frac{\partial^2V}{\partial x^2}=rV(t,x)
\end{aligned}
\]

and boundary condition $V(T,x)=f(x)$ for all $x$.

Under risk-neutral measure, $\dd S=rS\,\dd t+\sigma S\,\dd W(t)$. Let $S=X$, $\beta(t,X)=rS$ and
$\gamma(t,X)=\sigma S$, then the discounted Feynman-Kac equation becomes the
Black-Scholes-Merton differential equation

\[
\begin{aligned}
\frac{\partial V}{\partial t}+rS\frac{\partial V}{\partial S}\\
&\quad+\frac{1}{2}\sigma^2S^2\frac{\partial^2V}{\partial S^2}=rV.
\end{aligned}
\]

\end{solution}

\subsubsection*{Black-Scholes formula}

The Black-Scholes formula for European calls and puts with continuous dividend yield $y$
is:

\[
c=Se^{-y\tau}N(d_1)-Ke^{-r\tau}N(d_2)\quad\text{and}\quad
p=Ke^{-r\tau}N(-d_2)-Se^{-y\tau}N(-d_1),
\]

% ZHOU-SOURCE-PAGE: 160 PRINTED: 144

\[
\begin{aligned}
d_1&=\frac{\ln(Se^{-y\tau}/K)+(r+\sigma^2/2)\tau}{\sigma\sqrt{\tau}}\\
&=\frac{\ln(S/K)+(r-y+\sigma^2/2)\tau}{\sigma\sqrt{\tau}}
\end{aligned}
\]

where

\[
d_2=\frac{\ln(S/K)+(r-y-\sigma^2/2)\tau}{\sigma\sqrt{\tau}}=d_1-\sigma\sqrt{\tau}.
\]

$N(x)$ is the $\CDF$ of the standard normal distribution and $N'(x)$ is the $\PDF$ of the
standard normal distribution:

\[
N(x)=\int_{-\infty}^{x}\frac{1}{\sqrt{2\pi}}e^{-y^2/2}\,\dd y
\quad\text{and}\quad N'(x)=\frac{1}{\sqrt{2\pi}}e^{-x^2/2}.
\]

If the underlying asset is a futures contract, then yield $y=r$. If the underlying asset is a
foreign currency, then yield $y=r_c$, where $r_c$ is the foreign risk-free interest rate.

\begin{problem}{A. What are the assumptions behind the Black-Scholes formula?}
\end{problem}

\begin{solution}
The original Black-Scholes formula for European calls and puts consists of the equations
$c=SN(d_1)-Ke^{-r\tau}N(d_2)$ and
$p=Ke^{-r\tau}N(-d_2)-SN(-d_1)$, which require the following assumptions:

\begin{enumerate}
\item The stock pays no dividends.
\item The risk-free interest rate is constant and known.
\item The stock price follows a geometric Brownian motion with constant drift $\mu$ and
volatility $\sigma$: $\dd S=\mu S\,\dd t+\sigma S\,\dd W(t)$.
\item There are no transaction costs or taxes; the proceeds of short selling can be fully
invested.
\item All securities are perfectly divisible.
\item There are no risk-free arbitrage opportunities.
\end{enumerate}

\end{solution}

\begin{problem}{B. How can you derive the Black-Scholes formula for a European call on a non-dividend
paying stock using risk-neutral probability measure?}
\end{problem}

\begin{solution}
The Black-Scholes formula for a European call on a non-dividend paying stock is

\[
\begin{aligned}
c&=SN(d_1)-Ke^{-r\tau}N(d_2),\quad\text{where}\quad\\
d_1&=\frac{\ln(S/K)+(r+\sigma^2/2)\tau}{\sigma\sqrt{\tau}}\\
&\quad\text{and}\quad d_2=d_1-\sigma\sqrt{\tau}.
\end{aligned}
\]

% ZHOU-SOURCE-PAGE: 161 PRINTED: 145

Under the risk-neutral probability measure, the drift of stock price becomes the risk-free
interest rate $r(t)$: $\dd S=r(t)S\,\dd t+\sigma S\,\dd W(t)$. Risk-neutral measure allows the option to be
priced as the discounted value of its expected payoff with the risk-free interest rate:

\[
\begin{aligned}
V(t)&=\E\left[ e^{-\int_t^T r(u)\,\dd u}V(T)\vert S(t)\right],\\
&\quad 0\le t\le T,\quad\text{where }V(T)\text{ is the payoff at maturity }T.
\end{aligned}
\]

When $r$ is constant, the formula can be further simplified as
$V(t)=e^{-r\tau}\E[V(T)\vert S(t)]$.
Under risk-neutral probabilities, $\dd S=rS\,\dd t+\sigma S\,\dd W(t)$. Applying Ito's lemma, we get

\[
\begin{aligned}
\dd(\ln(S))&=(r-\sigma^2/2)\,\dd t+\sigma\,\dd W(t)\\
&\Rightarrow \ln S_T\sim N\left(\ln S+(r-\sigma^2/2)\tau,\sigma^2\tau\right).
\end{aligned}
\]

So $S_T=Se^{(r-\sigma^2/2)\tau+\sigma\sqrt{\tau}\epsilon}$, where $\epsilon\sim N(0,1)$. For a
European option, we have

\[
V(T)=
\begin{cases}
Se^{(r-\sigma^2/2)\tau+\sigma\sqrt{\tau}\epsilon}-K,
&\text{if }Se^{(r-\sigma^2/2)\tau+\sigma\sqrt{\tau}\epsilon}>K,\\
0,&\text{otherwise}.
\end{cases}
\]

\[
\begin{aligned}
Se^{(r-\sigma^2/2)\tau+\sigma\sqrt{\tau}\epsilon}>K
&\Rightarrow \epsilon>\frac{\ln(K/S)-(r-\sigma^2/2)\tau}{\sigma\sqrt{\tau}}=-d_2\\
&\quad\text{and}
\end{aligned}
\]

\[
\begin{aligned}
\E[V(T)\vert S]&=\E[\max(S_T-K,0)\vert S]\\
&=\int_{-d_2}^{\infty}\left(Se^{(r-\sigma^2/2)\tau+\sigma\sqrt{\tau}\epsilon}-K\right)
\frac{1}{\sqrt{2\pi}}e^{-\epsilon^2/2}\,\dd\epsilon\\
&=Se^{r\tau}\int_{-d_2}^{\infty}\frac{1}{\sqrt{2\pi}}
e^{-(\epsilon-\sigma\sqrt{\tau})^2/2}\,\dd\epsilon
-K\int_{-d_2}^{\infty}\frac{1}{\sqrt{2\pi}}e^{-\epsilon^2/2}\,\dd\epsilon.
\end{aligned}
\]

Let $\tilde{\epsilon}=\epsilon-\sigma\sqrt{\tau}$, then
$\dd\epsilon=\dd\tilde{\epsilon}$, $\epsilon=-d_2\Rightarrow
\tilde{\epsilon}=-d_2-\sigma\sqrt{\tau}=-d_1$ and we have

\[
\begin{aligned}
Se^{r\tau}\int_{-d_2}^{\infty}\frac{1}{\sqrt{2\pi}}
e^{-(\epsilon-\sigma\sqrt{\tau})^2/2}\,\dd\epsilon
&=Se^{r\tau}\int_{-d_1}^{\infty}\frac{1}{\sqrt{2\pi}}e^{-\tilde{\epsilon}^2/2}\,\dd\tilde{\epsilon}\\
&=Se^{r\tau}N(d_1),\\
K\int_{-d_2}^{\infty}\frac{1}{\sqrt{2\pi}}e^{-\epsilon^2/2}\,\dd\epsilon
&=K\left(1-N(-d_2)\right)=KN(d_2).
\end{aligned}
\]

\[
\begin{aligned}
\therefore\quad \E[V(T)]&=Se^{r\tau}N(d_1)-KN(d_2)\quad\text{and}\quad\\
V(t)&=e^{-r\tau}\E[V(T)]=SN(d_1)-Ke^{-r\tau}N(d_2).
\end{aligned}
\]

From the derivation process, it is also obvious that $1-N(-d_2)=N(d_2)$ is the
risk-neutral probability that the call option finishes in the money.

\end{solution}

\begin{problem}{C. How do you derive the Black-Scholes formula for a European call option on a
non-dividend paying stock by solving the Black-Scholes-Merton differential equation?}
\end{problem}

% ZHOU-SOURCE-PAGE: 162 PRINTED: 146

\begin{solution}
You can skip this problem if you don't have background in partial differential equations
(PDE). One approach to solving the problem is to convert the Black-Scholes-Merton
differential equation to a heat equation and then apply the boundary conditions to the heat
equation to derive the Black-Scholes formula.

Let $y=\ln S$ ($S=e^y$) and $\tilde{\tau}=T-t$, then
$\dfrac{\partial V}{\partial t}=-\dfrac{\partial V}{\partial\tilde{\tau}}$,
$\dfrac{\partial V}{\partial S}=\dfrac{\partial V}{\partial y}\dfrac{\dd y}{\dd S}
=\dfrac{1}{S}\dfrac{\partial V}{\partial y}$ and

\[
\begin{aligned}
\frac{\partial^2V}{\partial S^2}
&=\frac{\partial}{\partial S}\left(\frac{\partial V}{\partial S}\right)
=\frac{\partial}{\partial S}\left(\frac{1}{S}\frac{\partial V}{\partial y}\right)\\
&=-\frac{1}{S^2}\frac{\partial V}{\partial y}
+\frac{1}{S}\frac{\partial}{\partial S}\left(\frac{\partial V}{\partial y}\right)
=-\frac{1}{S^2}\frac{\partial V}{\partial y}
+\frac{1}{S^2}\frac{\partial^2V}{\partial y^2}.
\end{aligned}\footnotemark[6]
\]

\footnotetext[6]{The log is taken to convert the geometric Brownian motion to an arithmetic
Brownian motion; $\tau=T-t$ is used to convert the equation from a backward equation to a
forward equation with initial condition at $\tau=0$ (the boundary condition at $t=T\Rightarrow
\tau=0$).}

The Black-Scholes-Merton differential equation

\[
\begin{aligned}
\frac{\partial V}{\partial t}+rS\frac{\partial V}{\partial S}
&\quad+\frac{1}{2}\sigma^2S^2\frac{\partial^2V}{\partial S^2}-rV=0
\end{aligned}
\]

can be converted to

\[
\begin{aligned}
-\frac{\partial V}{\partial\tilde{\tau}}
&+\left(r-\frac{1}{2}\sigma^2\right)\frac{\partial V}{\partial y}\\
&\quad+\frac{1}{2}\sigma^2\frac{\partial^2V}{\partial y^2}-rV=0.
\end{aligned}
\]

Let $u=e^{r\tilde{\tau}}V$, the equation becomes

\[
\begin{aligned}
-\frac{\partial u}{\partial\tilde{\tau}}
&+\left(r-\frac{1}{2}\sigma^2\right)\frac{\partial u}{\partial y}\\
&\quad+\frac{1}{2}\sigma^2\frac{\partial^2u}{\partial y^2}=0.
\end{aligned}
\]

Finally, let
$x=y+\left(r-\dfrac{1}{2}\sigma^2\right)\tilde{\tau}
=\ln S+\left(r-\dfrac{1}{2}\sigma^2\right)\tilde{\tau}$ and $\tau=\tilde{\tau}$, then
$\dfrac{\partial u}{\partial y}=\dfrac{\partial u}{\partial x}$ and

\[
\begin{aligned}
\frac{\partial u}{\partial\tilde{\tau}}
&=\frac{\partial u}{\partial\tau}+\left(r-\frac{1}{2}\sigma^2\right)\frac{\partial u}{\partial x},
\end{aligned}
\]

which transforms the equation to

\[
\begin{aligned}
-\frac{\partial u}{\partial\tau}
-\left(r-\frac{1}{2}\sigma^2\right)\frac{\partial u}{\partial x}\\
&\quad+\left(r-\frac{1}{2}\sigma^2\right)\frac{\partial u}{\partial x}
+\frac{1}{2}\sigma^2\frac{\partial^2u}{\partial x^2}=0\\
&\Rightarrow\frac{\partial u}{\partial\tau}=\frac{1}{2}\sigma^2\frac{\partial^2u}{\partial x^2}.
\end{aligned}
\]

So the original equation becomes a heat/diffusion equation
$\dfrac{\partial u}{\partial\tau}=\dfrac{1}{2}\sigma^2\dfrac{\partial^2u}{\partial x^2}$. For heat
equation $\dfrac{\partial u}{\partial\tau}=\dfrac{1}{2}\sigma^2\dfrac{\partial^2u}{\partial x^2}$,
where $u=u(x,\tau)$ is a function of time $\tau$ and space variable $x$, with boundary
condition $u(x,0)=u_0(x)$, the solution is

\[
\begin{aligned}
u(x,\tau)&=\frac{1}{\sqrt{2\pi\sigma^2\tau}}\\
&\quad\int_{-\infty}^{\infty}u_0(\psi)\exp\left(-\frac{(x-\psi)^2}{2\sigma^2\tau}\right)\,\dd\psi.\footnotemark[7]
\end{aligned}
\]
